\subsubsection{Euler tour / Hierholzer (camino y circuito)}

\textbf{Teoria de Grafos Eulerianos ($O(V + E)$)}:
Recorre cada \textbf{arista} exactamente una vez. Todas las aristas deben pertenecer a una sola componente conexa (se ignoran vertices aislados de grado 0).
\begin{itemize}\setlength\itemsep{0pt}
	\item \textbf{Circuito Euleriano}: Empieza y termina en el \textbf{mismo} vertice ($A \to \dots \to A$).
	\item \textbf{Camino Euleriano}: Empieza en $A$ y termina en $B$ con $B \ne A$ ($A \to \dots \to B$).
	\item \textbf{No dirigido - Circuito}: Todos los vertices tienen \textbf{grado par}.
	\item \textbf{No dirigido - Camino}: Exactamente \textbf{dos vertices tienen grado impar} (inicio $A$ y fin $B$); todos los demas tienen grado par.
	\item \textbf{Dirigido - Circuito}: Para todo vertice $v$, $\text{in}(v) = \text{out}(v)$.
	\item \textbf{Dirigido - Camino}: Inicio $A$ con $\text{out}(A) - \text{in}(A) = 1$, fin $B$ con $\text{in}(B) - \text{out}(B) = 1$, y para todos los demas $\text{in}(v) = \text{out}(v)$.
	\item \textbf{Algoritmo de Hierholzer}: DFS que consume aristas; en post-order se apila el vertice y se invierte al final ($O(V + E)$).
\end{itemize}

\begin{lstlisting}[style=mystyle, language=C++]
// TC: O(N + M) - Algoritmo de Hierholzer
// usa: existencia y construccion de camino / circuito euleriano
//     (visita cada ARISTA exactamente una vez; admite multigrafos)
#include <bits/stdc++.h>
using namespace std;

// ---------- Condiciones necesarias y suficientes ----------
//   NO DIRIGIDO, conexo:
//     - circuito euleriano: todos los vertices tienen grado par
//     - camino  euleriano: exactamente 0 o 2 vertices con grado impar
//                            (si son 2, el camino va de uno al otro)
//   DIRIGIDO, debilmente conexo sobre el subgrafo con aristas:
//     - circuito euleriano: indeg(v) == outdeg(v) para todo v
//     - camino  euleriano: existe s con outdeg(s) = indeg(s)+1,
//                          existe t con indeg(t) = outdeg(t)+1,
//                          y el resto cumple indeg == outdeg

// ---------- Existencia (no dirigido) ----------
// returns: 0 = circuito, 1 = camino (start!=end), -1 = no existe
int eulerianStatusUndirected(int n, const vector<vector<int>>& adj){
	int odd = 0, any = -1;
	for(int v=0;v<n;v++){
		int d = (int)adj[v].size();
		if(d & 1) odd++;
		if(d > 0 && any == -1) any = v;
	}
	if(any == -1) return 0; // grafo vacio -> circuito trivial
	// chequeo rapido de conexion entre vertices con aristas
	vector<int> vis(n, 0);
	queue<int> q; q.push(any); vis[any] = 1;
	while(!q.empty()){
		int v = q.front(); q.pop();
		for(int u : adj[v]) if(!vis[u]){ vis[u] = 1; q.push(u); }
	}
	for(int v=0;v<n;v++) if(adj[v].size() && !vis[v]) return -1;
	if(odd == 0) return 0;
	if(odd == 2) return 1;
	return -1;
}

// ---------- Existencia (dirigido) ----------
// devuelve: 0 = circuito, 1 = camino, -1 = no existe; *start = nodo inicial
int eulerianStatusDirected(int n, const vector<vector<int>>& adj,
int& start){
	int moreOut = -1, moreIn = -1, bad = 0;
	start = 0;
	for(int v=0;v<n;v++){
		if(!adj[v].empty()) start = v;
		int indeg = 0;
		for(int u=0;u<n;u++) for(int x : adj[u]) if(x == v) indeg++;
		int outdeg = (int)adj[v].size();
		int diff = outdeg - indeg;
		if(diff == 1){ if(moreOut != -1) bad = 1; moreOut = v; }
		else if(diff == -1){ if(moreIn != -1) bad = 1; moreIn = v; }
		else if(diff != 0) bad = 1;
	}
	if(bad) return -1;
	if(moreOut == -1) return 0;      // circuito
	if(moreIn == -1) return -1;
	start = moreOut;
	return 1;                        // camino
}

// ---------- Hierholzer (no dirigido, multigrafo) ----------
// adj guarda (vecino, id) para cada arista; cada arista no dirigida aparece
//     en adj[u] y adj[v]. Tras llamar, adj queda mutado con ids 0..M-1.
// Retorna la secuencia de vertices; si es circuito, path.front() == path.back().
vector<int> eulerianPathUndirected(int n, vector<vector<pair<int,int>>>& adj){
	// reconstruye el grafo asignando ids consistentes
	vector<vector<pair<int,int>>> g(n);
	int idx = 0;
	for(int u=0;u<n;u++){
		for(auto& [v, idOld] : adj[u]){
			if(u < v){ // procesa cada arista no dirigida una sola vez
				g[u].push_back({v, idx});
				g[v].push_back({u, idx});
				idx++;
			}
		}
	}
	vector<int> vis(idx, 0); // vis de aristas (eid), no de vertices
	adj = move(g);

	// elegir inicio segun existencia (trabajamos con la lista de vecinos)
	vector<vector<int>> nbrs(n);
	for(int u=0;u<n;u++) for(auto& [v, eid] : adj[u]) nbrs[u].push_back(v);
	int start = 0, status = eulerianStatusUndirected(n, nbrs);
	if(status == 1){
		for(int v=0;v<n;v++) if(((int)adj[v].size()) & 1){ start = v; break; }
	}else if(status == 0){
		for(int v=0;v<n;v++) if(!adj[v].empty()){ start = v; break; }
	}else return {};

	stack<int> st;
	vector<int> path, it(n, 0);
	st.push(start);
	while(!st.empty()){
		int v = st.top();
		while(it[v] < (int)adj[v].size() && vis[adj[v][it[v]].second]) it[v]++;
		if(it[v] < (int)adj[v].size()){
			auto [u, eid] = adj[v][it[v]];
			vis[eid] = 1;
			st.push(u);
		}else{
			path.push_back(v);
			st.pop();
		}
	}
	reverse(path.begin(), path.end());
	// si el camino no usa todas las aristas, no existe euleriano
	int total = 0; for(int x : vis) total += x;
	if(total != (int)vis.size()) return {};
	return path;
}

// ---------- Hierholzer (dirigido) ----------
vector<int> eulerianPathDirected(int n, vector<vector<pair<int,int>>>& adj){
	// reescribir ids consistentes
	vector<vector<pair<int,int>>> g(n);
	int idx = 0;
	for(int u=0;u<n;u++){
		for(auto& [v, idOld] : adj[u]){
			g[u].push_back({v, idx++});
		}
	}
	vector<int> vis(idx, 0); // vis de aristas (eid), no de vertices
	adj = move(g);

	// existencia (sobre la lista de vecinos)
	vector<vector<int>> nbrs(n);
	for(int u=0;u<n;u++) for(auto& [v, eid] : adj[u]) nbrs[u].push_back(v);
	int start = 0;
	int status = eulerianStatusDirected(n, nbrs, start);
	if(status == -1) return {};

	stack<int> st;
	vector<int> path, it(n, 0);
	st.push(start);
	while(!st.empty()){
		int v = st.top();
		while(it[v] < (int)adj[v].size() && vis[adj[v][it[v]].second]) it[v]++;
		if(it[v] < (int)adj[v].size()){
			auto [u, eid] = adj[v][it[v]];
			vis[eid] = 1;
			st.push(u);
		}else{
			path.push_back(v);
			st.pop();
		}
	}
	reverse(path.begin(), path.end());
	int total = 0; for(int x : vis) total += x;
	if(total != (int)vis.size()) return {};
	return path;
}

int main(){
	// ----- no dirigido: triangulo 0-1-2-0 tiene circuito euleriano (todos grado par) -----
	vector<vector<int>> adjU(3);
	adjU[0] = {1,2}; adjU[1] = {0,2}; adjU[2] = {0,1};
	int stU = eulerianStatusUndirected(3, adjU);
	cout << "no dirigido status: " << stU << " (0=circuito, 1=camino)\n";

	// ----- dirigido: 0->1->2->0 es circuito; agregamos 0->3->2 para tener camino 0->2 -----
	vector<vector<int>> adjD(4);
	adjD[0] = {1,3}; adjD[1] = {2}; adjD[2] = {0}; adjD[3] = {2};
	int startD = 0;
	int stD = eulerianStatusDirected(4, adjD, startD);
	cout << "dirigido status: " << stD << " start=" << startD << "\n";

	// reconstruir un circuito euleriano en el triangulo
	vector<vector<pair<int,int>>> agU(3);
	agU[0].push_back({1, 0}); agU[1].push_back({0, 0});
	agU[1].push_back({2, 1}); agU[2].push_back({1, 1});
	agU[2].push_back({0, 2}); agU[0].push_back({2, 2});
	auto pathU = eulerianPathUndirected(3, agU);
	for(int v : pathU) cout << v << " ";
	cout << "\n";
	return 0;
}
\end{lstlisting}
